\section[电磁对偶]{电磁对偶（2 月 4 日）}

\srcpage{18}

上次我们讨论了 Cartier 对偶，即 Pontrjagin 对偶的代数几何版本，例如用 $\mathbb{G}_m$（或者对总在 $\C$ 上工作的各位来说，用 $\C^\times$）代替 $U_1$。不过代数故事与分析故事确实能互相提供有用的东西。

$G$ 的表示论有不同版本，它们对应对偶世界中不同的群代数概念。例如，若把群代数看成 von Neumann 代数，对偶 $\widehat{G}$ 具有测度空间的结构。若考虑群 $C^*$ 代数，$\widehat{G}$ 是拓扑空间。而若只考虑离散代数（于是 $G$ 很小），则 $\widehat{G}$ 是代数簇。你应当把这些看成一个对偶对象 $\widehat{G}$，尽管精确的数学对象并不相同。

现在来看一些例子，尽管它们其实都是同一个例子——$\Z$ 与圆之间的 Fourier 级数——只是说得更花哨一些。

\begin{example}[Pontrjagin--Poincaré 对偶]
  设 $M$ 是紧 $n$ 维流形，$G$ 是局部紧阿贝尔群；于是可用 $M$ 作为阿贝尔群的来源，具体来说是 $H_i(M;G)$。Pontrjagin--Poincaré 对偶指的是存在完美配对
  \[
    H_i(M;G)\otimes H^{n-i}(M;\widehat{G})\longrightarrow H^n(M;G\otimes\widehat{G})\longrightarrow H^n(M;U_1)\longrightarrow U_1,
  \]
  从而自然地把 $H^{n-i}(M;\widehat{G})$ 等同于 $H_i(M;G)$ 的 Pontrjagin 对偶。在同伦论中，这与一个相当一般的、抽象版本的 Pontrjagin 对偶有关，称为 Brown--Comenetz 对偶。

  （TODO：如果 $G$ 带拓扑，会发生什么？我想确保自己弄明白这一点。）

  特殊化到 $1$ 维，我们得到从 $\Z$（或映射 $\mathrm{pt}\to\Z$）到 $\C^\times$（或映射 $\mathrm{pt}\to\C^\times$）的对偶。这不太令人兴奋。但可以把它看成 $\Z$ 上的量子力学与 $\C^\times$ 上的量子力学之间的等价。
\end{example}

\srcpage{19}

那 $2$ 维呢？现在可以考虑 $H^0(S^1;\Z)$，我们已知它对偶于 $H^1(S^1;\C^\times)$。特别地，有一个 Fourier 变换把这些阿贝尔群上的函数等同起来。可以把 $H^0(S^1;\Z)$ 看成 $[S^1,\Z]$，即到 $\Z$ 的映射的同伦类。对 $H^1(S^1;\C^\times)$，我们想称它为 $\mathrm{Loc}_{\C^\times}(S^1)$，但这本应指叠的版本；所以就让 $\mathrm{Loc}_{\C^\times}(S^1)$ 表示其底空间吧。

来说说局部系统。$G$-局部系统是一个局部同构于常值层 $G$ 的层，且其粘接映射等同于 $G$ 通过乘法作用在自身上。某种意义上，这并不用到 $G$ 的拓扑或几何，只用到底层集合。例如，局部系统有单值映射（monodromy）$\pi_1(X)\to G$，并且它在同构意义下由这个映射决定。换个说法，局部系统的同构类由 $H^1(X;G)$ 给出——即使 $G$ 非阿贝尔也如此！通常我们只习惯对阿贝尔群定义上同调；但对非阿贝尔群，可以在 $0$ 与 $1$ 次定义奇异上同调或 Čech 上同调，而 $H^1$ 恢复了主丛的同构类。

局部系统有自同构，记住这一点往往很重要；但现在我们只看同构类。当我们只看底空间 $\mathrm{Loc}_{\C^\times}(S^1)$ 时，它就是空间 $\C^\times$，于是又得到 $\Z$ 与 $\C^\times$ 之间的 Cartier 对偶。这看起来无聊，但既然我们升了一维、见到了主丛，这个对偶就与 $U_1$ 的规范场论有关了。

\begin{example}[另一个二维例子]
  现在考虑 $H^1(S^1;\Z)$，我们称之为 A 侧。其对偶 $H^0(S^1;\C^\times)$ 称为 B 侧。这可不是在说唱片的正反面！

  在 B 侧，我们看的是到 $\C^\times$ 的映射，其中 $\C^\times$ 取离散拓扑。在物理中这是标量场的一个例子。在 A 侧，我们看的是 $[S^1,B\Z]$；由于 $B\Z = K(\Z,1) = S^1$，我们看的是到 $S^1$ 的映射。这是另一个标量场，你可以把它看成 T-对偶或镜像对称的一个极特殊情形。在 T-对偶中，这两个圆的半径互换 $R\leftrightarrow 1/R$。

  若想把这件事做得更精致些，可以把 $\Z$ 换成格 $\Lambda$，这样就没有典范基，对偶也更有趣一点。此时对偶是对偶复环面 $T^\vee_\C$；而 $H^1(S^1;\Lambda) = [S^1,T]$，其中 $T$ 是紧环面 $T := \Lambda\otimes_\Z S^1$。这又是 T-对偶的一个例子。
\end{example}

\begin{example}
  现在升到 $3$ 维。设 $\Sigma$ 是紧致定向 $2$-流形。Pontrjagin--Poincaré 对偶把 $H^1(\Sigma;\Z)$ 与 $H^1(\Sigma;\C^\times)$ 互换。左边是标量场 $[\Sigma,B\Z]$，或 $[\Sigma,S^1]$；右边得到 $\mathrm{Loc}_{\C^\times}(\Sigma)$，这关乎 $3$ 维规范场论。

  另一种做法是把 $H^2(\Sigma;\Z)$ 与 $H^0(\Sigma;\C^\times)$ 互换。此时 $H^2(\Sigma;\Z)\cong H^1(\Sigma;U_1) = [\Sigma,BU_1]$，因为 $BU_1 = K(\Z,2)$。这关乎主 $U_1$-丛的规范场论。右边得到 $[\Sigma,\C^\times]$，即标量场的理论。同样，若想更有趣，可以升级到环面与格，于是得到向量空间的同构。
\end{example}

\begin{example}
  现在是一个四维例子：$H^2(M;\Z)$ 与 $H^2(M;\C^\times)$。有挠时也能让这套东西跑起来；但为了叙述方便，假定 $H^2(M;\Z)$ 无挠。维数 $4$ 的妙处在于，这个对应两侧都能等同于规范场论：A 侧 $H^2(M;\Z) = [M,K(\Z,2)] = [M,BU_1]$，于是得到主 $U_1$-丛。换个说法，$[M,BU_1] = \pi_0(\mathrm{Map}(M,\C P^\infty))$。主 $U_1$-丛与复线丛之间有等价：一个方向取主 $U_1$-丛的相伴线丛，另一个方向是选一个度量再取单位圆丛。于是可以把 $H^2(M;\Z)$ 看成 $M$ 的 Picard 群——不过若你习惯的是带更多结构（例如簇）的 Picard 群，这可能是个奇怪的版本。

  B 侧得到 $\C^\times$-值局部系统，至少在不计叠性时如此。

  这是电磁对偶的一个实例。参见 Witten 在 quantum fields and strings 丛书中的讲座「Abelian duality I」与「Abelian duality II」，以及 Freed [Fre00] 与 Freed--Moore--Segal [FMS07] 中更数学化的处理。
\end{example}

（TODO：我漏掉了一些东西）我们想研究方程 $du = 0$ 与 $d\star u = 0$ 的解，其中 $u = d\varphi$。第一个条件自动成立，因为 $u$ 恰当。第二个条件等价于 $\star d\star u = 0$，或等价地，$\varphi$ 是调和的。

但对调和条件的这一改写显示出一个漂亮的对称性：$\{du = 0,\ d\star u = 0\}$ 在 Hodge 星下对称。由于 $\star u$ 闭，至少局部地有 $\star u = d\varphi^\vee$。于是 $U_1$-值标量场的理论对偶于另一个 $U_1$-值标量场的理论。或者，若考虑映入某个环面的映射，对偶侧就是映入对偶环面的映射。当你引入度量以便更好地谈论调和性时，会看到：若 A 侧圆的半径是 $R$，则 B 侧圆的半径是 $1/R$。

现在看看 $3$ 维里发生什么。设 $\varphi\colon\Sigma\to S^1$ 是光滑映射，$u = d\varphi$。施加 $du = 0$ 与 $d\star u = 0$，对偶于要求 $F = \star u = dA$。

\srcpage{20}

但 $4$ 维中的对称性更有意思。电磁学的速成课：有电场 $E\in\Omega^1(\R^3)$ 与磁场 $B\in\Omega^2(\R^3)$。若想在 Minkowski 空间上相对论式地工作，就把它们并成场强
\[
  F := B - dt\wedge E\in\Omega^2(\mathbb{M}^4),
\]
其中 $\mathbb{M}^4$ 表示 Minkowski 空间。

真空中的 Maxwell 方程告诉我们 $dF = 0$ 与 $d\star F = 0$。这当然在 $F\leftrightarrow\star F$ 下对称。若破坏相对论对称性并写
\[
  \star F = B^\vee - dt\wedge E^\vee,
\]
则 $B^\vee = -\star_3 E$、$E^\vee = \star_3 B$，其中 $\star_3$ 表示在时间切片（一个 $3$-流形）上取 Hodge 星。（也就是说，空间是 $3$ 维的，时空是 $4$ 维的。）

若加入带流的电磁电荷，方程会改变：此时有电流 $j_E$ 与磁流 $j_B$，它们是 $3$-形式，我们要求 $dF = j_B$、$d\star F = j_E$。在一个时间切片 $M$ 上，$\int_M j_B = Q_B$、$\int_M j_E = Q_E$，分别是磁荷与电荷。Stokes 定理给出 Gauss 定律：若 $\Sigma\subset M$ 闭，则磁通量
\[
  b_\Sigma = \int_\Sigma F = -Q_B.
\]
这可以解释为离开 $\Sigma$ 的场线的平均数。类似地，总电通量为
\[
  e_\Sigma = \int_\Sigma \star F = Q_E.
\]

（TODO：别以为 Arun 的符号都是对的！）

自然界没有磁单极子，所以 $Q_B = 0$；Gauss 定律说，若你在闭曲面上积分，得到 $0$。这允许我们引入一个新对象——电磁势 $A$，并写 $F = dA$。这破坏了电与磁之间的对称性。给定这些数据，$\nabla := d + A$ 是 $M$ 上主 $U_1$-丛的联络，而规范对称性
\[
  \begin{aligned}
    \nabla &\longmapsto g^{-1}\nabla g,\\
    A &\longmapsto A + g^{-1}dg,
  \end{aligned}
\]
其中 $g\colon M\to U_1$ 光滑，不影响 $F$。

这是个非常方便的招术，但它在物理上有意义吗？有！有一个实验上被测量过的原理，叫 Aharonov--Bohm 效应，它赋予 $A$ 物理意义。即使 $F = 0$（联络平坦），带电粒子沿环路也会获得相位，这测量了平坦联络的和乐。所以 $A$ 确实存在（尽管只在规范变换意义下）。

既然我们在 $3$-流形上把 $F$ 换成了 $dA$（并引入了 $A$），也来看看在 $4$ 维里这样做会怎样。这个选择现在实现了 Dirac 电荷量子化。设 $M = N\times\R$，其中 $N$ 是 $3$ 维的、$M$ 是 $4$ 维的；设 $\nabla$ 是主 $U_1$-丛 $P\to M$ 上的联络。即使没有带电粒子，我们仍有通量，因为可以积分
\[
  b_\Sigma = \frac{1}{2\pi i}\int_\Sigma F.
\]
它不再为零；事实上由 Chern--Weil 理论，它就是 $\langle c_1(P),[\Sigma]\rangle$。这是个整数，而不是任意的复数——这就是 Dirac 电荷量子化。电通量也有类似的东西：
\[
  e_\Sigma = \frac{1}{2\pi i}\int_\Sigma \star F.
\]

经典故事到此为止，接下来转向量子故事。我们将在 Hamilton 形式下讨论它。量子场论意味着：我们不是用点粒子做量子力学，而是用场来做。场是那些具有局部性、可以层化的东西（这不是精确定义！）：到某个空间的映射、向量丛的截面、主丛或向量丛连同联络，诸如此类。

\srcpage{21}

作为启发式，$M^{d-1}\times\R$ 上的量子场论应当感觉像 $M^{d-1}$ 上场空间上的量子力学。这个空间常常是无限维的。这很像分析中的标准手法：把偏微分方程看成无限维的常微分方程。

我们将有态构成的 Hilbert 空间 $\mathcal{H}$，启发式地它就是 $M$ 上场空间的 $L^2$。量子场论所做的部分事情，就是把 Hilbert 空间附着到时间切片上。眼下不管正则性问题：$L^2$ 的意思是「构成 Hilbert 空间的函数」，但我们注意到存在良好的解。

于是我们想考虑空间 $\mathcal{C}(M)$，其元素是 $M$ 上带酉联络的主线丛 $L\to M$ 的同构类。联络总有自同构，暂时忽略这一点。

所以我们尝试在 Hilbert 空间 $L^2(\mathcal{C}(M))$ 上做量子力学。$\mathcal{C}(M)$ 有很好的结构：在线丛与联络的张量积下，它是 Lie 群的无限维类比；并且某种意义上它分裂为一个格、一个环面与一个无限维向量空间的乘积。于是今天讨论的阿贝尔对偶的两侧都在这里，外加我们想研究 Fourier 级数的那个向量空间。

格来自一个离散不变量，即第一 Chern 类 $c_1\colon\mathcal{C}(M)\to H^2(M;\Z)$；在物理中它测量磁通量。几何上，也可以取联络的曲率 $F\colon\mathcal{C}(M)\to\Omega^2(M)$ 来恢复它。但它不是随便一个 $2$-形式：它闭，且有整周期（意思是它与 $H^2(M;\Z)$ 对话）。曲率映射的核是平坦联络构成的空间 $\mathcal{C}^\flat(M)$，它等同于 $M$ 上 $U_1$-值局部系统（即主 $U_1$-丛，其中 $U_1$ 取离散拓扑，并取这些丛的同构类）。我们讨论过，同构意义下的局部系统可以看成 $H^1(M;\R/\Z)$，其中认同 $\R/\Z\cong U_1$；而 $H^1(M;\R/\Z)$ 是环面 $U_1\otimes_\Z H^1(M;\Z)$，假定 $H^1(M;\Z)$ 无挠。

格 $\Lambda := H^2(M;\Z)$ 称为磁通量格；在 $\mathcal{C}(M)$ 内部有局部系统的环面 $T = \mathrm{Loc}_{U_1}(M)$，此外还有向量空间的部分。于是我们知道 $L^2$ 长什么样：对 $T$ 与 $\Lambda$ 有 Fourier 级数；对向量空间的 $L^2$，有 Fourier 变换——尽管它是无限维的。

从这个视角看，电磁对偶将在 $\mathcal{C}(M)$ 上应用 Pontrjagin 对偶，我们下次讨论。我们将能把同一个 Hilbert 空间写成两种样子；一侧的磁测量就是另一侧的电测量。
